\subsection{AI and the Concentration of Power in Tech}
\label{sec:8.3.2}
A small group of firms — Google, Apple, Meta, Amazon, and Microsoft — holds a share of AI development disproportionate to their number, and that concentration carries economic and ethical stakes independent of what any single one of them does with the position it occupies. No individual development team escapes this pull, for reasons already given.

Vast financial resources let firms of this size fund extensive research and development, draw the top talent in the field, and command sprawling data networks — the raw material effective machine learning depends on. That combination widens the gap between these firms and smaller companies or new entrants, who cannot match the data or the resources and therefore cannot compete on equal terms.

Acquisition sharpens the same dynamic. Facebook's acquisitions of Instagram and WhatsApp — the company now operates as Meta — did not just remove two competitors; it handed one firm access to comprehensive behavioral data from billions of users worldwide \autocite{ftc2020metaplatforms}. Consolidation of that kind extends a firm's control over how AI gets used in social media, advertising, and e-commerce, and from inside a market that concentrated the trajectory looks like a foregone conclusion.

Power concentrated this tightly is also power positioned for exploitation, and the risk sharpens considerably when the technology in question reaches an authoritarian government. The relationship between the Chinese state and firms like Alibaba and Tencent is the clearest live example: state access to those firms' machine learning capability has been documented as a channel for surveillance and censorship at scale, which is a serious risk to free expression, individual privacy, and democratic accountability wherever the same model gets exported \autocite{hrw2020repression,cnn2020alibabauyghur}.

Concentration has a quieter cost alongside the loud ones. A firm controlling the datasets and infrastructure a competitive system needs is a gatekeeper whether or not it intends to be, and fewer independent attempts at AI development means fewer chances at the kind of breakthrough an incumbent has no incentive to pursue itself.

The obvious countermeasure is antitrust, and it has been tried against exactly the acquisitions described above. The Federal Trade Commission sued Meta in December 2020 to unwind the Instagram and WhatsApp deals. After a six-week bench trial the district court held, in November 2025, that the agency had not shown Meta currently holds monopoly power — TikTok and YouTube compete for the same attention, so a market defined without them was too narrow — and the FTC appealed the following January \autocite{ftc2020metaplatforms}. Whatever the appeal decides, the case is a lesson about the instrument rather than about this defendant: a merger challenge brought a decade after the mergers has to prove a present-day market, and the market had moved. That leaves the remedies which do not depend on winning that argument — mandated access to the datasets that function as chokepoints, and public funding for research outside the dominant firms. Section~\ref{sec:6.4.3} develops export controls, prohibited-practices regulation and industry design standards for the narrower case of authoritarian misuse; the concern here is prior to that one: concentration itself, whatever any single firm does with the position it holds.
